\begin{frame}{스팸 필터: 임계값을 움직여보자 (설정)}
  \textbf{같은 모델}의 예측 확률을 임계값 0.5, 0.6, 0.1로 자를
  때 혼동행렬이 어떻게 달라지는지 --- ``모델 성능''은 사실
  \kb{``임계값 + 모델''}의 성능.
  \vskip0.6em
  \begin{itemize}
    \item \textbf{스팸 필터} 두 특징: $x_1$ = 스팸 단어 빈도
      (실제 메일에서 10$\sim$36개 범위), $x_2$ = ``발신자
      미등록'' 여부 (0/1).
    \item 정규 메일은 스팸 단어 적고 대부분 발신자 등록.
    \item 스팸의 70\%는 ``분명한'' 스팸, 30\%는 정규 메일과 겹치는
      \textbf{``교묘한'' 스팸} --- 임계값을 아무리 올려도 잡을 수
      없는 스팸이 항상 남음.
    \item 전체 500통 중 스팸 100통 (20\%), \textbf{테스트 세트
      150통} (스팸 23통, 15.3\%).
  \end{itemize}
  {\footnotesize \texttt{StandardScaler}로 표준화 후 본문
  \texttt{logistic\_gradient\_descent}로 학습 --- 스케일링 필수:
  $x_1$ 의 스케일(10$\sim$36)이 $x_2$ (0/1)의 30배라, 안 하면
  2.1절 ``학습률의 함정''이 그대로 살아남고 왜곡된 GD가 된다.}
\end{frame}
